% TOP_PROBLEM: {"id":30007540,"problem_number":"LOCAL-30007540","title":"Effective fiscal-rule design","category_id":21,"category":{"id":21,"name":"economics","display_name":"Economics","description":"Open economic theory statements, published research agendas, and proposed scoped specifications; question provenance and historical priority are qualified per record.","slug":"economics","order_index":21,"created_at":"2026-10-08T00:00:00Z"},"status":"provenance_pending","external_url":null,"legacy_ids":[],"published":false,"statement":"Design incentive-compatible fiscal rules and escape clauses when compliance, circumvention, and political benefits are endogenous.","economics":{"original_id":29,"field":"Fiscal policy and sovereign debt","kind":"design","formalization_basis":"proposed_specification","source_status":"not_verified","priority_status":"not_established","priority_note":"No explicit primary open-problem statement matching this catalogue item was verified; supporting research is not evidence of first formulation.","earliest_verified_reference":null,"current_reference":{"authors":["Julien Acalin","Virginia Alonso-Albarran","Clara Arroyo","Waikei Raphael Lam","Leonardo Martinez","Anh D. M. Nguyen","Francisco Roch","Galen Sher","Alexandra Solovyeva"],"title":"Fiscal Guardrails against High Debt and Looming Spending Pressures","year":2025,"url":"https://www.imf.org/en/publications/staff-discussion-notes/issues/2025/09/22/fiscal-guardrails-against-high-debt-and-looming-spending-pressures-569841","doi":"10.5089/9798229023801.006","locator":"Primary Staff Discussion Note 2025/004 page, “Main findings” and “Policy implications”","evidence_summary":"The report finds mixed compliance and discusses rule revisions, escape clauses, institutions, and correction mechanisms. This supports institutional motivation but does not state the exact optimal-rule problem below as open."},"formalization_note":"New optimization specification. The IMF report supports the compliance/escape-clause motivation but no precise primary open statement of this target was verified."},"statusQualification":"Provenance pending: an explicit published statement of this exact open question was not verified. This is a catalogue-authored candidate specification, not a certified open conjecture. New optimization specification. The IMF report supports the compliance/escape-clause motivation but no precise primary open statement of this target was verified.","statusReviewedAt":"2026-10-08"}
% FIRST_POSTED: 2026-10-08
% ENTRY_KIND: statement_only
% !TeX program = lualatex
\documentclass[11pt]{article}
\usepackage{fontspec}
\usepackage[a4paper,margin=25mm]{geometry}
\usepackage{amsmath,amssymb,mathtools}
\usepackage{xurl}
\usepackage[hidelinks]{hyperref}
\newcommand{\sref}[2]{\hyperlink{source:#1:#2}{[S#2]}}
\setlength{\emergencystretch}{3em}
\begin{document}
% problemId:30007540
\section*{Effective fiscal-rule design}\label{30007540}
\subsection{Definitions and mathematical statement}
Fix a stochastic fiscal economy with debt ratio \(b_t\), output growth \(g_t\), real financing rate \(r_t\), and primary surplus \(s_t\), expressed as a ratio to output at \(t+1\), with \(g_t>-1\). Debt evolves as \(b_{t+1}=(1+r_t)b_t/(1+g_t)-s_t\). A political government values resident welfare plus a declared electoral benefit from current spending. Fiscal institutions select rule parameters \(a=(\phi,\bar b,e,p)\): a debt-feedback coefficient, reference debt ratio, measurable escape-clause trigger, and enforcement probability or penalty schedule. The government retains explicitly bounded discretion when the trigger activates.
For a stated admissible rule family \(\mathcal A\), identify the rule maximizing expected resident welfare subject to budget feasibility, government incentive compatibility, and a bound on crisis risk. Formally,
\[
a^*\in\arg\max_{a\in\mathcal A}E_a\sum_{t=0}^{H}\beta^t u(c_t,\ell_t).
\]
Allow compliance, attempted circumvention, and rule revisions to be endogenous rather than assuming the announced rule is always followed. Estimate the response of those behaviors to monitoring and enforcement using credible institutional variation, and quantify the insurance value of escape clauses against their opportunistic use. Compare alternative rule designs under common shocks and political preferences. This is a proposed institutional-design problem; successful empirical associations between rules and debt performance do not establish an optimal universally applicable rule.

\par\textbf{Formulation and scope.} New optimization specification. The IMF report supports the compliance/escape-clause motivation but no precise primary open statement of this target was verified.

\par\textbf{Open-question provenance.} An explicit published open-question statement was not verified. Sources below are supporting literature and must not be treated as a first listing.

\par\textbf{Historical priority.} No explicit primary open-problem statement matching this catalogue item was verified; supporting research is not evidence of first formulation.

\subsection{Short English statement}
Design incentive-compatible fiscal rules and escape clauses when compliance, circumvention, and political benefits are endogenous.

\subsection{Sources}
\begin{itemize}\raggedright
\item[\textnormal{[S1]}]\hypertarget{source:30007540:1}{} Julien Acalin, Virginia Alonso-Albarran, Clara Arroyo, Waikei Raphael Lam, Leonardo Martinez, Anh D. M. Nguyen, Francisco Roch, Galen Sher, Alexandra Solovyeva. 2025. Fiscal Guardrails against High Debt and Looming Spending Pressures. Primary Staff Discussion Note 2025/004 page, “Main findings” and “Policy implications”.
\url{https://www.imf.org/en/publications/staff-discussion-notes/issues/2025/09/22/fiscal-guardrails-against-high-debt-and-looming-spending-pressures-569841}.
DOI: 10.5089/9798229023801.006.
{\small\textit{Supporting or current literature; see evidence qualification. The report finds mixed compliance and discusses rule revisions, escape clauses, institutions, and correction mechanisms. This supports institutional motivation but does not state the exact optimal-rule problem below as open.}}
\end{itemize}
\end{document}
